\section*{Dédicaces}

Il existe des réussites que l’on signe seule, mais que l’on ne construit jamais vraiment sans les autres.

Aujourd’hui, alors que s’achève une étape importante de ma vie, je souhaite dédier ce travail à celles et ceux qui ont été ma force dans les moments de doute, mon refuge dans les moments difficiles et ma joie dans les instants heureux. Derrière ces pages, il y a bien plus que des heures de travail : il y a votre amour, votre patience et toutes ces petites attentions qui m’ont permis de continuer à avancer.

\textbf{À mon père, Ramzi Khouni, et à ma mère, Nedia Rezigui}, merci pour votre amour inconditionnel, vos sacrifices, votre patience et cette confiance que vous m’avez accordée, même lorsque je peinais à croire en moi-même. Derrière chaque étape franchie se trouvent vos encouragements, vos efforts souvent silencieux et votre présence rassurante. Si je suis arrivée jusqu’ici, c’est aussi parce que vous m’avez appris à ne pas abandonner, même lorsque le chemin devient difficile. Ce diplôme est le fruit de mon travail, mais il est aussi, et surtout, une partie du vôtre. \textbf{J’espère qu’en regardant le chemin parcouru, vous ressentirez ne serait-ce qu’une partie de la fierté que j’éprouve à être votre fille.}

\textbf{À mes grands-parents, Yamina, Sihem et Belgacem}, merci pour votre tendresse, vos prières, votre sagesse et cet amour si particulier que seuls les grands-parents savent offrir. Votre présence est une richesse que les mots ne suffisent pas à exprimer. Vous représentez ces racines qui nous donnent la force de grandir sans jamais oublier d’où nous venons.

J’adresse une pensée particulièrement émue à \textbf{mon grand-père Habib, que Dieu ait son âme en paix}. Certaines absences se font ressentir plus profondément lors des moments que l’on aurait tant aimé partager. Même s’il n’est plus parmi nous pour assister à cette étape de ma vie, sa mémoire et l’affection que je lui porte continuent de m’accompagner. Il y a des personnes que le temps nous enlève, mais que le cœur ne cesse jamais de garder près de lui. \textbf{Ce moment est aussi pour toi, grand-père.}

\textbf{À ma tante Hedia Khouni et à mon oncle Sifi}, j’adresse mes pensées les plus chaleureuses. À toi, Hedia, qui m’as toujours entourée d’un amour et d’une attention dignes d’une seconde mère, merci de ne m’avoir jamais simplement considérée comme ta nièce, mais comme ta propre fille. Ta générosité, ta tendresse, ta présence et tes encouragements ont été pour moi une source précieuse de réconfort et de force. Il existe des liens que les mots « tante » et « nièce » ne suffisent pas à décrire, et le nôtre en fait partie. À toi, Sifi, merci pour ta gentillesse, ta disponibilité et ton soutien tout au long de mon parcours. Je vous suis sincèrement reconnaissante à tous les deux pour votre présence et pour la place particulière que vous occupez dans ma vie.

\textbf{À mes chères cousines, Jouhaina et Yosra Khedhiri}, merci pour votre complicité, votre affection et tous ces souvenirs qui rendent la vie plus belle. Merci pour nos fous rires, nos conversations interminables, nos petites folies et ces moments où, sans même le chercher, vous réussissez à me redonner le sourire. Avec vous, même les instants les plus ordinaires deviennent des souvenirs précieux. Vous êtes bien plus que des cousines, vous êtes des amies, des confidentes et une partie de mon bonheur. Et même si je ne vous le dis pas toujours (il ne faudrait quand même pas trop vous flatter !), je suis infiniment reconnaissante de vous avoir dans ma vie.

\textbf{À vous tous, ma famille}, merci d’avoir supporté mes inquiétudes, mes longues journées de travail et mes explications techniques parfois interminables. Vous n’avez peut-être pas toujours compris pourquoi un simple bug pouvait bouleverser toute ma journée, mais vous avez toujours su trouver les mots, les gestes ou les sourires pour me redonner du courage.

À toutes les personnes qui ont croisé mon chemin, qui m’ont tendu la main, offert un conseil, partagé un sourire ou simplement encouragée au bon moment, merci. Il arrive que quelques mots suffisent à redonner l’élan nécessaire pour continuer, et ces gestes, même discrets, ne s’oublient pas.

\textbf{Si mon nom figure sur la couverture de ce rapport, une grande partie de la force qui m’a permis de l’écrire vient de vous. Et si je devais retenir une seule réussite au-delà de ce diplôme, ce serait la chance de vous avoir dans ma vie.}

Ce travail, je vous le dédie avec tout mon amour et toute ma reconnaissance.

\textbf{Parce qu’au fond, la plus belle des réussites n’est pas seulement d’arriver au bout du chemin, mais de pouvoir partager ce moment avec les personnes qui lui donnent tout son sens.}


\section*{Remerciements}

Il est difficile de résumer en quelques lignes une aventure qui m’a tant appris, parfois même bien au-delà de ce que j’avais imaginé en la commençant.

Ce projet de fin d’études marque l’aboutissement de plusieurs années de travail, de découvertes, de remises en question et de défis. Il représente bien plus qu’une étape académique : c’est une expérience humaine, faite de rencontres précieuses, de petites victoires, de grandes interrogations et de quelques imprévus qui, avec le recul, mériteraient presque leur propre chapitre.

Avant toute chose, je tiens à exprimer ma profonde gratitude à \textbf{Monsieur Mohamed Yassine Ben Debba}, gérant d’\textbf{IPROTEX Tunisie}, pour la confiance qu’il m’a accordée, son accueil et l’intérêt qu’il a porté à mon projet. Je le remercie sincèrement de m’avoir offert l’opportunité de découvrir le monde industriel de l’intérieur et de transformer une idée en une solution répondant à des besoins concrets. Cette confiance a été pour moi une véritable source de motivation et de responsabilité. J’adresse également mes sincères remerciements à \textbf{Madame Dorsaf Ben Debba}, pour sa bienveillance, sa disponibilité et son soutien tout au long de cette expérience.

Je souhaite exprimer toute ma reconnaissance à \textbf{Madame Sarra Slimani}, mon encadrante académique, pour ses conseils, son accompagnement et ses remarques constructives. Ses exigences m’ont poussée à aller plus loin, à questionner mes choix et à ne pas me contenter de la première solution qui fonctionne. Grâce à son regard critique, j’ai appris qu’en ingénierie, terminer quelque chose est une satisfaction, mais comprendre pourquoi et comment l’améliorer en est une autre.

Mes remerciements s’adressent également à \textbf{l’ensemble des collaborateurs d’IPROTEX}, qui ont pris le temps de partager leurs connaissances, de répondre à mes questions et de m’aider à comprendre un univers industriel qui m’était, au départ, peu familier. Leur patience a été précieuse, surtout face à ma curiosité qui semblait parfois avoir autant de questions qu’une machine possède de composants. Grâce à eux, j’ai compris qu’avant de chercher à digitaliser un processus, il faut apprendre à écouter celles et ceux qui le connaissent et le font vivre au quotidien.

Je remercie chaleureusement \textbf{l’ensemble des enseignants d’ESPRIT}, qui ont contribué à construire les fondations de mon parcours. Au fil de ces années, ils m’ont transmis bien plus que des connaissances : une méthode de travail, un esprit critique et cette capacité à chercher des solutions, même lorsque le problème semble décidé à ne pas coopérer.

Une place toute particulière dans ces remerciements revient à \textbf{l’IEEE ESPRIT Student Branch} et à \textbf{l’IEEE ESPRIT SIGHT Group}. Ces communautés ont été bien plus que des engagements associatifs : elles ont été des lieux de rencontres, d’apprentissage, d’amitiés et de découvertes. J’y ai appris à prendre des initiatives, à travailler avec des personnes aux idées différentes, à assumer des responsabilités et, parfois, à improviser lorsque les choses ne se déroulaient absolument pas comme prévu. Elles m’ont surtout appris que l’ingénierie ne consiste pas uniquement à concevoir des systèmes performants, mais aussi à mettre ses compétences au service des autres. Certaines expériences se terminent avec un diplôme ; celles-ci laisseront des souvenirs, des valeurs et des liens que j’espère conserver bien plus longtemps.

Enfin, je voudrais accorder quelques lignes au chemin parcouru. À ces moments où tout semblait avancer, puis à ceux où un simple détail suffisait à remettre en question des heures de travail. Aux idées qui ont évolué, aux erreurs qui ont enseigné davantage que certaines réussites, et à cette persévérance que j’ai parfois découverte en moi sans vraiment savoir qu’elle s’y trouvait.

Ce projet m’a appris que l’on ne maîtrise pas toujours les circonstances, que les solutions les plus intéressantes naissent parfois des difficultés les plus inattendues et qu’en informatique, même après avoir tout vérifié, il reste toujours une petite chance que le problème disparaisse mystérieusement après un redémarrage.

Mais surtout, il m’a appris une chose essentielle : \textbf{la réussite ne se résume pas à ce que l’on accomplit, mais aussi à ce que l’on surmonte, aux personnes qui nous accompagnent et à celle que l’on devient en chemin.}

À toutes les personnes qui ont contribué à cette aventure, qui m’ont accordé leur confiance, partagé leurs connaissances ou simplement encouragée à avancer, j’adresse ma plus sincère reconnaissance.

\textbf{À vous tous, du fond du cœur, merci.}
